\documentclass[11pt,a4paper]{article}
\usepackage[T1]{fontenc}
\usepackage[utf8]{inputenc}
\usepackage{lmodern}
\usepackage{amsmath,amssymb,amsthm,mathtools}
\usepackage{esint,graphicx}
\usepackage[a4paper,margin=22mm,headheight=14pt]{geometry}
\usepackage{microtype}
\usepackage{booktabs,longtable,array,enumitem}
\newcolumntype{L}[1]{>{\raggedright\arraybackslash}p{#1}}
\usepackage[dvipsnames]{xcolor}
\usepackage{fancyhdr}
\usepackage{hyperref}

\newtheorem{theorem}{Theorem}[section]
\newtheorem{lemma}[theorem]{Lemma}
\newtheorem{proposition}[theorem]{Proposition}
\newtheorem{corollary}[theorem]{Corollary}
\theoremstyle{remark}
\newtheorem{remark}[theorem]{Remark}
\newtheorem{definition}[theorem]{Definition}
\newtheorem{claim}[theorem]{Claim}
\setlist{nosep,leftmargin=1.5em}
\setlength{\emergencystretch}{2em}
\pagestyle{fancy}
\fancyhf{}
\fancyhead[L]{\small AIM 679}
\fancyhead[R]{\small Revised 7 November 2026}
\fancyfoot[C]{\thepage}
\renewcommand{\headrulewidth}{0.3pt}
\hypersetup{colorlinks=true,linkcolor=MidnightBlue,urlcolor=MidnightBlue,citecolor=MidnightBlue,pdftitle={AIM 679},pdfauthor={Matthew J. Colbrook; results found by AI}}
\begin{document}
\title{The sequence \texorpdfstring{$\sqrt{k}/k!$}{sqrt(k)/k!} is a multiplier sequence}
\author{Matthew J. Colbrook\thanks{The solutions and partial results in this document were found by AI. The supplied research and the separate repository review are AI work; human peer review and formal verification are not claimed.}\\[4pt]
\small Department of Applied Mathematics and Theoretical Physics\\
\small University of Cambridge, Cambridge CB3 0WA, United Kingdom\\
\small\href{mailto:m.colbrook@damtp.cam.ac.uk}{m.colbrook@damtp.cam.ac.uk}}
\date{7 November 2026}
\maketitle
\begin{abstract}
$\{\sqrt{k}/k!\}_{k\ge0}$ is a multiplier sequence; its exponential generating function has simple negative nonzero zeros with explicit quadratic localization. The conclusion concerns exactly the catalogue target; no historical priority claim is made.
\end{abstract}
\noindent\textbf{Catalogue scope.} AIM 679; complete resolution. \href{https://github.com/MColbrook/AIM/blob/1e564633494c109edfd3ec46da1366adcff925a5/problems/679-the-square-root-factorial-multiplier-sequence.md}{Canonical target at revision \texttt{1e56463}}.


\section{Statement and argument}
\label{sec:679}

The answer to the question of Csordas and Forg\'acs
\cite{CsordasForgacs} is affirmative.  We prove the stronger assertion
that the associated entire function has only simple negative zeros,
apart from its simple zero at the origin, and locate every zero between
two explicit bounds.

\begin{theorem}
The sequence $\gamma_0=0$, $\gamma_k=\sqrt{k}/k!$ for $k\ge1$, is a
classical multiplier sequence.  More precisely, write
\[
 \Phi(w)=\sum_{k=1}^{\infty}\frac{\sqrt{k}}{(k!)^2}w^k.
\]
Its zeros are $0$ and $-\lambda_j$, $j\ge1$, all simple, where
\[
 \frac{j^2\pi^2}{4}<\lambda_j<\frac{(j+1)^2\pi^2}{4}.
\]
In particular $\Phi$ belongs to the Laguerre--P\'olya class of type I,
with nonnegative Taylor coefficients.
\end{theorem}

\subsection{A cosine subtraction lemma}

\begin{lemma}
Let $\mu$ be a nonatomic probability measure on $(0,1)$.  The entire
function
\[
 F(z)=\cos z-\int_{(0,1)}\cos(tz)\,d\mu(t)
\]
has only real zeros.  Its zero at $0$ has multiplicity two; in each
interval $(j\pi,(j+1)\pi)$, $j\ge1$, it has exactly one simple zero.
The remaining zeros are the negatives of these zeros.
\end{lemma}

\begin{proof}
Local uniform domination proves that the integral is entire.  The
function $F$ is even and has real Taylor coefficients.  Since
$0<\int t^2\,d\mu(t)<1$,
\[
 F(z)=-\frac{1-\int t^2\,d\mu(t)}{2}z^2+O(z^4),
\]
so the zero at $0$ is exactly double.

For a nonzero integer $j$ and $z=j\pi+iy$,
\[
 |\cos z|=\cosh y,
 \qquad
 \left|\int\cos(tz)\,d\mu(t)\right|
       \le\int\cosh(ty)\,d\mu(t).
\]
When $y\ne0$ the second quantity is strictly less than $\cosh y$,
because $0<t<1$.  When $y=0$ use instead
\[
 \left|\int\cos(j\pi t)\,d\mu(t)\right|
 \le\int|\cos(j\pi t)|\,d\mu(t)<1.
\]
Here the equality set of $|\cos(j\pi t)|=1$ in $(0,1)$ is finite and
has zero $\mu$-measure.

On a horizontal line $z=x\pm iR$, uniformly in real $x$,
\[
 |\cos z|\ge\sinh R,
 \qquad
 \left|\int\cos(tz)\,d\mu(t)\right|
       \le\int\cosh(tR)\,d\mu(t).
\]
Moreover
\[
 \frac{\int\cosh(tR)\,d\mu(t)}{\sinh R}\longrightarrow0
       \quad(R\longrightarrow\infty).
\]
Indeed, the integrand ratio tends to zero for every $t<1$ and is bounded
by $\coth1$ for $R\ge1$.  Dominated convergence applies.

For every sufficiently large $R$, Rouch\'e's theorem on the rectangle
with vertical sides $j\pi$ and $(j+1)\pi$, $j\ge1$, and horizontal
sides at imaginary parts $\pm R$ therefore gives exactly one zero of
$F$, counted with multiplicity: this is the zero count for $\cos z$.
Conjugation invariance forces this zero to be real, and it is simple.
Letting $R$ increase excludes any other zero in the strip.  On the
central rectangle with vertical sides $-\pi$ and $\pi$, the same
argument gives two zeros counted with multiplicity, already accounted
for by the double zero at $0$.  Evenness handles the negative strips.
The strict vertical-side estimates exclude zeros on their boundaries.
These regions cover the plane.
\end{proof}

\subsection{A moment representation for a gamma ratio}

\begin{lemma}
For $x>0$ define
\[
 W(x)=\sqrt{x}\,\frac{\Gamma(x+\tfrac12)}{\Gamma(x+1)}.
\]
There is a nonatomic probability measure $\mu$ on $(0,1)$ such that
\[
 W(x)=1-\int_{(0,1)}t^{2x}\,d\mu(t),\qquad x>0.
\]
\end{lemma}

\begin{proof}
For $s>0$ put
\begin{align*}
 K(s)
 &=\frac1\pi\int_0^s\frac{dt}{\sqrt{s-t}\sqrt{e^t-1}}\\
 &=\frac1\pi\int_0^1\frac1{\sqrt{1-u}}
                   \sqrt{\frac{s}{e^{su}-1}}\,du.
\end{align*}
Since $e^{su}-1\ge su$, the integrand is bounded by the integrable
function $\pi^{-1}[u(1-u)]^{-1/2}$.  For fixed $0<u<1$,
$s/(e^{su}-1)$ is strictly decreasing in $s$: the sign of its derivative
is the sign of $e^v(1-v)-1<0$, where $v=su>0$.  Thus $K$ is continuous
and strictly decreasing.  Dominated convergence gives
\[
 K(0+)=\frac1\pi\int_0^1[u(1-u)]^{-1/2}\,du=1,
 \qquad K(\infty)=0.
\]
It follows that $K$ is the survival function of a nonatomic probability
measure $\nu$ on $(0,\infty)$, so that
$\nu((s,\infty))=K(s)$.

Tonelli's theorem, the gamma integral, and the beta integral show that
\begin{align*}
 \int_0^\infty e^{-xs}K(s)\,ds
 &=\frac1\pi
   \left(\int_0^\infty e^{-xr}r^{-1/2}\,dr\right)
   \left(\int_0^\infty e^{-xt}(e^t-1)^{-1/2}\,dt\right)\\
 &=\frac1\pi\frac{\sqrt\pi}{\sqrt x}
                 B\left(x+\tfrac12,\tfrac12\right)\\
 &=\frac1{\sqrt x}\frac{\Gamma(x+\tfrac12)}{\Gamma(x+1)}.
\end{align*}
For the second integral the substitution $q=e^{-t}$ gives
$\int_0^1q^{x-1/2}(1-q)^{-1/2}\,dq$.
The survival-function identity, also directly obtained by Tonelli, is
\[
 1-\int_0^\infty e^{-xs}\,d\nu(s)
  =x\int_0^\infty e^{-xs}K(s)\,ds=W(x).
\]
Push $\nu$ forward by $s\mapsto e^{-s/2}$.  The resulting measure $\mu$
is nonatomic, has total mass one, is supported in $(0,1)$, and has the
required moment identity.
\end{proof}

\subsection{Identification and factorization}

Let $\mu$ be the measure just constructed.  Termwise integration of the
locally uniformly convergent cosine series and the gamma duplication
formula give
\begin{align*}
 \frac1{\sqrt\pi}
 \left(\cos z-\int\cos(tz)\,d\mu(t)\right)
 &=\frac1{\sqrt\pi}\sum_{k=1}^\infty
            \frac{(-1)^kW(k)}{(2k)!}z^{2k}\\
 &=\sum_{k=1}^\infty\frac{\sqrt{k}}{(k!)^2}
                         \left(-\frac{z^2}{4}\right)^k
 =\Phi\left(-\frac{z^2}{4}\right).
\end{align*}
The coefficient identity here is
\[
 \frac{\Gamma(k+\tfrac12)}{\sqrt\pi\,\Gamma(k+1)}
       =\frac{(2k)!}{4^k(k!)^2}.
\]
The cosine subtraction lemma now implies that all zeros of $\Phi$ are
nonpositive real numbers.  In fact every complex number $w$ has a
square root $z$ with $w=-z^2/4$, and if $\Phi(w)=0$ then the lemma
forces this $z$ to be real.  Multiplicities and the interval bounds in
the theorem follow from the same substitution.

For completeness, $\Phi$ has order at most $1/2$.  For $r\ge0$,
\[
 M_\Phi(r)\le\sum_{k\ge1}\frac{k r^k}{(k!)^2}
   =\sum_{k\ge1}\frac{r^k}{(k-1)!k!}
   \le\sqrt r\,e^{2\sqrt r}.
\]
The last inequality follows by comparison with the product of two
exponential series.  Hadamard factorization for an entire function of
order less than one therefore gives
\[
 \Phi(w)=w\prod_{j=1}^\infty\left(1+\frac{w}{\lambda_j}\right).
\]
The factor in front is exactly $w$ because $\Phi'(0)=1$.  The product
converges locally uniformly, also immediately from
$\lambda_j>j^2\pi^2/4$.  Its finite products have only nonpositive
zeros and nonnegative coefficients.  Hence $\Phi$ is in the required
Laguerre--P\'olya class.  The P\'olya--Schur characterization, in the
normalization $\Phi(w)=\sum\gamma_k w^k/k!$, proves that
$\{\sqrt{k}/k!\}_{k\ge0}$ is a multiplier sequence.

\begingroup\small
\begin{thebibliography}{99}
\bibitem{CsordasForgacs} G. Csordas and T. Forgacs, \emph{Multiplier sequences, classes of generalized Bessel functions and open problems}, J. Math. Anal. Appl. \textbf{433} (2016), 1369--1389. \url{https://arxiv.org/abs/1505.00355v2}. DOI: \url{https://doi.org/10.1016/j.jmaa.2015.08.047}.
\end{thebibliography}
\endgroup
\end{document}
